\documentclass[11pt]{article}
\usepackage[margin=0.9in]{geometry}
\usepackage{amsmath,amssymb,amsthm,booktabs,mathrsfs}
\usepackage[T1]{fontenc}
\usepackage{lmodern,microtype}
\usepackage[colorlinks=true,linkcolor=blue,urlcolor=blue]{hyperref}
\newtheorem{theorem}{Theorem}[section]
\newtheorem{lemma}[theorem]{Lemma}
\newtheorem{proposition}[theorem]{Proposition}
\newtheorem{corollary}[theorem]{Corollary}
\theoremstyle{remark}
\newtheorem{remark}[theorem]{Remark}
\newcommand{\norm}[1]{\left\|#1\right\|}
\title{Toward a four-dimensional Yang--Mills construction\\
\large Notebook 17: joint fluctuation tails and the one-step Gaussian domain}
\author{Research working document}
\date{17 October 2026}
\begin{document}
\maketitle
\begin{abstract}
We prove joint bad-block bounds for correlated Gaussian fluctuations
from a counting-coordinate covariance bound. The estimates allow
singular observable covariances and overlapping tests. A relative
complex quadratic perturbation has a corresponding local estimate
without a factor growing with the total volume: untested variables
are integrated out before taking absolute values. In four dimensions,
an exact fixed-ratio rescaling supplies the required covariance bound
for both potential values and first differences. The actual reduced
Wilson Hessian at our nonlinear constrained minimizer satisfies the
relative quadratic hypothesis, including its constraint-curvature
term. A separate extraction of the earlier flat effective-action
bounds supplies the current-lattice Gaussian at arbitrary homogeneous
averaging depth. These are inputs to the quantum integration and
iteration required for the existence argument.
\end{abstract}

\section{Which fluctuation estimate is needed}
\label{gs17-target}

Our objective is a nontrivial four-dimensional Yang--Mills quantum
theory. Notebooks 15 and 16 provide the homogeneous finite-regulator
classical minimizer, its ordinary-constraint transfer, local analytic
responses, and the minimized action. The next step integrates the
fluctuations and sums the regions that leave the analytic chart.

The Gaussian roughness calculation of Notebook 16 distinguishes two
different integrations. Integrating every unresolved mode below a
fixed macroscopic scale does not give a uniform smooth-field domain.
Integrating a fixed-ratio step has different coordinates and different
uniform bounds. Both statements must retain the physical pairing:
an inverse precision as a physical Hilbert operator is not its raw
coordinate covariance matrix.

This notebook proves the joint tail estimate needed for prescribed
bad blocks, its complex version, and the fixed-step quadratic
identifications supported by our construction. A good-cell estimate
alone does not supply a regional nonlinear chart on its complement.
The latter construction and the interacting weights are explicit
interfaces for the next step.

The proofs have been checked internally by separate derivations and
audits. The results are finite-regulator theorems on their stated
domains; a continuum existence theorem is the target of their
subsequent integration.


% No preamble. One-step quantum interface. Prefix os-.
\section{One fixed-ratio fluctuation step in physical coordinates}
\label{os-section}

We work in dimension four with a fixed odd blocking ratio
$L\ge3$. Fine spacing is $\eta>0$ and the next retained
spacing is $h=L\eta$. The number of retained cells may be
arbitrarily large; constants may depend on $L$ and the gauge
group, but not on $h$, $\eta$, or compatible periodic volume.
This is one blocking step, rather than a simultaneous
integration of all fluctuations below a fixed macroscopic scale.

Let $y=x/h$, and write the rescaled potential as
$A_s(y)=hA_{\rm phys}(x)$. Then
\begin{equation}\label{os-rescaling}
 U(e)=e^{i\eta A_{\rm phys}(e)}=e^{iA_s(e)/L},\qquad
 \nabla_\eta A_{\rm phys}=h^{-2}\nabla_{1/L}A_s.
\end{equation}
The four-dimensional coupling-free Wilson action is unchanged
by this relabeling: it is the same sum of plaquette densities,
each counted once. The physical Maxwell action likewise obeys
\begin{equation}\label{os-quadratic-scaling}
 \eta^4\sum|d_\eta A_{\rm phys}|^2
       =L^{-4}\sum|d_{1/L}A_s|^2.
\end{equation}
The corresponding scalar gauge projection and flat averaging
constraint are the fixed-$L$ unit-coarse-scale operators of
the preceding notebooks after this rescaling.

Write $X_L=\ker Q_L\cap\ker(R_LD_*^T)$, where
$D_*=D_{1/L}$, and let $\mathcal C_L$ be the Maxwell
variational inverse on $X_L$. For the physical constrained
Gaussian $Z_{\rm phys}$ and its rescaling $Z_s=hZ_{\rm phys}$,
the covariance \emph{operators} under the physical pairings
and their raw coordinate matrices satisfy
\begin{equation}\label{os-covariance-scaling}
 \mathcal C_h=h^2\mathcal C_L,\qquad
 \operatorname{Cov}_{\rm count}(Z_{\rm phys})
       =\eta^{-4}h^2\mathcal C_L,\qquad
 \operatorname{Cov}_{\rm count}(Z_s)=L^4\mathcal C_L.
\end{equation}
The last two are coordinate covariances, rather than operators
under weighted physical sums. They follow also by pushing the
normalized Gaussian through \eqref{os-quadratic-scaling};
the constant Lebesgue Jacobian cancels from its normalization.

Define the complete dimensionless value/first-difference vector
\begin{equation}\label{os-microscopic-field}
 Y=\bigl(\eta Z_{\rm phys},\eta^2\nabla_\eta Z_{\rm phys}\bigr)
   =\bigl(Z_s/L,\nabla_*Z_s/L^2\bigr).
\end{equation}
Equation \eqref{os-covariance-scaling} gives the exact counting
covariance
\begin{equation}\label{os-Y-covariance}
 \Sigma_Y=S_L\mathcal C_LS_L^T,
       \qquad S_L=(LI,\nabla_*).
\end{equation}
The previously proved energy sandwiches imply
\begin{equation}\label{os-uniform-covariance}
 0\le\Sigma_Y\le K_L I
\end{equation}
for a fixed finite $K_L$, uniformly in spacing and volume.
Indeed $\mathcal C_L\le G_{0,L}$ and
$(I,\nabla_*)G_{0,L}(I,\nabla_*)^T$ is uniformly bounded;
the extra factor $L$ is fixed. All real trace-orthonormal
color components can be included at once.

Every retained cell and a fixed collar contain a fixed number
$m_L$ of these scalar tests. A test belongs to at most a fixed
number $\nu_L$ of such collars. Consequently a covariance-based
joint-tail theorem applies with constants independent of the
physical scale and the number of cells. Singularity of
$\Sigma_Y$ from gradients or constraints causes no problem for
such an upper bound.

\subsection{Small nonconstant backgrounds: the actual reduced Hessian}
\label{os-nonflat}

We can also control the quadratic fluctuation form at the
nonlinear corrected minimizer from Notebook 15, on its stated
small global coarse-data ball. Everything in this subsection
is first expressed at retained spacing one, with fine spacing
$1/L$. Let
\[
 A=\xi+H[B-\mathcal R(A)],\qquad
 A_*=A(\xi_*(B),B),\qquad
 T_A=(I+HD\mathcal R(A))^{-1}.
\]
Use the physical trace pairing and let $K_X$ be the compression
of $K=d^Td$ to $X_L$. Write $\mathcal B_B$ for the Hessian,
in the linear $\xi$ coordinate, of the exact reduced Wilson
action at $\xi_*(B)$. This is a fine fluctuation Hessian;
it is different from the coarse Hessian $D_B^2V$ of Notebook 16.

For $z,w\in X_L$, put $v=T_{A_*}z$ and $u=T_{A_*}w$.
The exact second derivative of the constraint chart gives
\begin{equation}\label{os-reduced-form}
 \mathcal B_B[z,w]
  =\langle v,K u\rangle
   +\langle S_{\rm jet}v,W(A_*)S_{\rm jet}u\rangle
   -\langle m(B),D^2\mathcal F(A_*)[v,u]\rangle.
\end{equation}
Here $W(A_*)=Df_{1/L}(S_{\rm jet}A_*)$ has norm
$O(\norm{B}_\infty)$ and $m(B)=O(\norm{B}_\infty)$.
The last term is retained; a Hessian of the ambient action
alone would not be the reduced fluctuation form.

The previously established tangent estimate is
\begin{equation}\label{os-tangent-energy}
 \norm v_2+\norm{\nabla_*v}_2
       \le C\langle z,Kz\rangle_{\rm Herm}^{1/2},
\end{equation}
and the same holds for $u,w$, also on the common complex
background ball. To see its complex extension directly, use
$v=z-HD\mathcal R(A_*)v$, the physical Hilbert smallness
$\norm{D\mathcal R(A_*)}_{2\to2}=O(\norm B_\infty)$,
the bounded $H,\nabla_*H$, and the flat coercive bound on $X_L$.

The flat mixed terms vanish, giving the exact bilinear identity
\begin{equation}\label{os-K-difference}
 \langle v,Ku\rangle-\langle z,Kw\rangle
   =\langle D\mathcal R(A_*)v,
                      q D\mathcal R(A_*)u\rangle.
\end{equation}
The right side is $O(\norm B_\infty^2)$ relative to the two
Maxwell energy norms. The middle term of
\eqref{os-reduced-form} is $O(\norm B_\infty)$ by
\eqref{os-tangent-energy}. The genuine two-Hilbert-direction
constraint estimate bounds the last term by the same order.
Consequently, on a fixed smaller complex ball,
\begin{equation}\label{os-relative-bound}
 \mathcal B_B=K_X^{1/2}(I+E_B)K_X^{1/2},\qquad
 E_B^T=E_B,\qquad \norm{E_B}_{2\to2}\le C\norm B_\infty.
\end{equation}
The normalized matrix $E_B$ is holomorphic in $B$. The square
root is that of the fixed real positive $K_X$; the norm is
the Hermitian operator norm, while the displayed transpose
is the analytic bilinear transpose.

For real $B$, choose the radius so that $\norm{E_B}\le\epsilon<1$.
The centered Gaussian $\zeta_s\in X_L$ for this reduced
quadratic form has covariance operator
$\mathcal C_B=J_X\mathcal B_B^{-1}J_X^T$ and satisfies
\begin{equation}\label{os-real-comparison}
 (1+\epsilon)^{-1}\mathcal C_L
     \le\mathcal C_B\le(1-\epsilon)^{-1}\mathcal C_L.
\end{equation}
Its dimensionless \emph{chart-coordinate} vector
$(\zeta_s/L,\nabla_*\zeta_s/L^2)$ therefore has counting
covariance bounded by $(1-\epsilon)^{-1}K_LI$.

The linearized physical-potential fluctuation is
$T_{A_*}\zeta_s$. It also has a uniformly bounded
dimensionless value/gradient covariance. Indeed
\eqref{os-tangent-energy} bounds
$S_LT_{A_*}J_XK_X^{-1/2}$ uniformly, so
\begin{equation}\label{os-tangent-covariance}
 \norm{S_LT_{A_*}\mathcal C_BT_{A_*}^TS_L^T}_{2\to2}
       \le C_L(1-\epsilon)^{-1}.
\end{equation}
Here the finite-dimensional maps on $X_L$ are read with the
same inclusions as in \eqref{os-real-comparison}. This is a
real Gaussian statement at a real background.

For complex $B$, \eqref{os-relative-bound} instead defines a
normalized complex Gaussian functional with strictly positive
Hermitian real part of its precision. Estimates for fixed real
coordinate tests can be obtained from its finite-dimensional
marginals. A cutoff whose defining tests themselves depend on
complex $B$ is a different problem; no holomorphy is inferred
by placing an indicator around such moving complex tests.

\subsection{The local domain for the nonlinear chart}
\label{os-chart-domain}

At the fixed mesh $1/L$, bounded values control all physical
first differences and their H\"older seminorms with constants
depending only on $L$ and $\alpha$. Thus on the event that
\emph{all fine value coordinates in the regulator} of the
dimensionless chart vector have modulus at most $R$,
\begin{equation}\label{os-domain-radius}
 \norm{g\zeta_s}_{N_\alpha}\le C_{L,\alpha}gR.
\end{equation}
Choose $gR$ small enough to stay inside the nonlinear chart
and its analytic action domain. If $R(g)=(\log(1/g))^p$
for fixed $p>0$, then both $gR(g)$ and $gR(g)^3$ tend to
zero. This is compatible with a large Gaussian-tail threshold.
The same norm comparison on a fixed cell and its stencil
collar is a local comparison of fields. It does not extend
the globally defined nonlinear minimizer chart through an
arbitrary complement of bad cells; that requires a regional
construction with its own boundary estimates.
It is the relevant fixed-ratio local scaling, rather than
a uniform smooth bound on all unresolved macroscopic modes.

On arbitrarily large volumes the event that \emph{every} cell
is good need not have large probability. The quantum step
requires joint bounds on prescribed bad cells and a summation
over their connected regions. The covariance bound supplies
the first of these ingredients. Integrating the nonlinear
remainder, the Haar and coordinate Jacobians, and bad regions
remains a further task; the present Gaussian comparison does
not replace those integrals.

% No preamble. Notebook 17 joint-tail module. Unique labels: jt-.
% Requires amsmath, amssymb, amsthm and a one-argument \norm macro.

\section{Joint Gaussian bounds for prescribed bad blocks}
\label{jt-section}

The input to this section is a counting-coordinate covariance
operator bound. A bound on each coordinate variance alone does
not suffice. We first give a real Gaussian theorem, including
degenerate covariances and overlapping tested sets. We then prove
a normalized complex Gaussian extension under an explicit
relative quadratic hypothesis. All statements are initially
finite-dimensional; their constants do not depend on the ambient
dimension or the number of tested blocks.

\subsection{The covariance input and microscopic blocks}
\label{jt-input}

Let $Y=(Y_i)_{i\in I}$ be a centered real Gaussian vector with
counting-coordinate covariance
\begin{equation}\label{jt-covariance}
 0\le\Sigma\le K I,\qquad K>0.
\end{equation}
For each block $b$ let $I_b\subset I$ be a nonempty tested set.
Assume fixed constants $m,\nu$ satisfy
\begin{equation}\label{jt-incidence}
 |I_b|\le m,\qquad
 \sup_i\#\{b:i\in I_b\}\le\nu.
\end{equation}
Thus $m$ is the number of scalar tests per block and $\nu$ is
their overlap multiplicity. Define
\begin{equation}\label{jt-bad-event}
 \mathsf B_b(R)=\{\max_{i\in I_b}|Y_i|>R\},\qquad R>0.
\end{equation}
Tests may be strongly correlated. The overlap counts repeated
$Y$ coordinates, not overlap of the underlying field variables
from which a derivative test was formed. Enlarging a block to
include a fixed stencil changes $m$ and $\nu$ by fixed factors.

\begin{theorem}[Joint real Gaussian bad-block bound]
\label{jt-real-theorem}
For every finite set $S$ of $k$ blocks,
\begin{equation}\label{jt-real-bound}
 \mathbb P\!\left(\bigcap_{b\in S}\mathsf B_b(R)\right)
 \le\left[
 m\sqrt2\exp\!\left(-\frac{R^2}{4\nu K}\right)
 \right]^k.
\end{equation}
The prefactor $m\sqrt2$ may instead be $2^{m/2}$, and the
right side may be replaced by its minimum with one.
\end{theorem}

\begin{proof}
Choose one coordinate $i_b\in I_b$ for each block. The vector
$W=(Y_{i_b})_{b\in S}$ has covariance at most $\nu K I$:
the coordinate-selection map, with repetitions allowed, has
squared operator norm at most $\nu$. Diagonalizing this
nonnegative covariance, allowing zero eigenvalues, gives
\begin{equation}\label{jt-square-mgf}
 \mathbb E e^{t|W|^2}
 =\det(I-2t\Sigma_W)^{-1/2}
 \le(1-2t\nu K)^{-k/2},\qquad 0\le t<(2\nu K)^{-1}.
\end{equation}
If every chosen coordinate exceeds $R$ in absolute value,
$|W|^2>kR^2$. Take $t=(4\nu K)^{-1}$ in Markov's inequality.
The resulting bound is $2^{k/2}e^{-kR^2/(4\nu K)}$.
There are at most $m^k$ coordinate choices, and their union
contains the event in \eqref{jt-real-bound}.

Alternatively collect every tested coordinate in the selected
blocks, retaining repetitions. This vector has at most $mk$
entries and covariance at most $\nu K I$. On the all-bad event
its squared norm is at least $kR^2$. Applying
\eqref{jt-square-mgf} with dimension $mk$ proves the alternative
prefactor. Neither proof assumes independent blocks.
\end{proof}

If a local norm obeys
$N_b(Y)\le c_b\max_{i\in I_b}|Y_i|$, its bad event is
bounded by the theorem with threshold $R/c_b$. For a Euclidean
block norm, the second proof applies directly to
$\sum_{i\in I_b}Y_i^2>R^2$, with prefactor $2^{m/2}$.
For $gY$ replace $R$ by $R/g$, equivalently $K$ by $g^2K$.

\subsection{Exponential squares, entropy, and connected sets}
\label{jt-entropy-section}

Put $E_b(Y)=\sum_{i\in I_b}Y_i^2$. For nonnegative $t_b$ of
finite support define
\[
 a_i=\sum_{b:i\in I_b}t_b,\qquad A=\operatorname{diag}(a_i).
\]
Then $\norm A\le\nu t_{\max}$ and
$\operatorname{Tr}A\le m\sum_b t_b$. The exact Gaussian
determinant, applied to $\Sigma^{1/2}A\Sigma^{1/2}$, gives
\begin{equation}\label{jt-exponential-square}
 \boxed{\quad
 \log\mathbb E\exp\!\left(\sum_b t_bE_b(Y)\right)
 \le\frac{Km\sum_b t_b}{1-2K\nu t_{\max}},
 \qquad 2K\nu t_{\max}<1.
 \quad}
\end{equation}
Indeed each eigenvalue is at most $K\nu t_{\max}$, their
sum is at most $K\operatorname{Tr}A$, and
$-\log(1-x)\le x/(1-x_{\max})$ for $0\le x\le x_{\max}<1$.
This proof also permits singular covariance.

For a probability law $\rho$ absolutely continuous with respect
to the Gaussian law $\mu$, the entropy variational inequality
therefore gives
\begin{equation}\label{jt-relative-entropy}
 \sum_b t_b\mathbb E_\rho E_b(Y)
 \le D(\rho\Vert\mu)
       +\frac{Km\sum_b t_b}{1-2K\nu t_{\max}}.
\end{equation}
Truncation gives the same assertion with extended expectations
when necessary.

Let $p(R)$ be the bracket in \eqref{jt-real-bound}, decreased
to at most one if desired. Expanding a finite product with
$z_b\ge0$ yields
\begin{equation}\label{jt-bad-count-product}
 \mathbb E\prod_b(1+z_b\mathbf1_{\mathsf B_b(R)})
 \le\prod_b(1+p(R)z_b).
\end{equation}
In particular, for $N$ tested blocks and $s\ge0$, the
exponential moment of their bad-block count is at most
$(1+p(R)(e^s-1))^N$. This is a moment statement, not an
assertion that the process is independent or stochastically
dominated for every increasing event.

Give the blocks an adjacency graph of maximum degree at most
$\Delta\ge1$. The number of connected $k$-block sets containing
a specified root is at most $(4\Delta)^{k-1}$. One way to see
this is to choose a canonical rooted traversal tree: ordered
rooted tree shapes cost at most $4^{k-1}$, and each edge has
at most $\Delta$ neighbor choices. Overcounting only enlarges
the bound. Consequently, whenever $4\Delta e^a p(R)<1$,
\begin{equation}\label{jt-connected-sum}
 \boxed{\quad
 \mathbb E\sum_{\substack{S\ni b_0\\S\ \mathrm{connected}}}
 e^{a|S|}\prod_{b\in S}\mathbf1_{\mathsf B_b(R)}
 \le\frac{e^a p(R)}{1-4\Delta e^a p(R)}.
 \quad}
\end{equation}
This is a summable Gaussian bad-set estimate. It does not
estimate an interacting polymer activity or prove its strict
background dependence.

\subsection{The homogeneous flat applications}
\label{jt-flat-applications}

For the previously constructed d4 constrained flat Gaussian,
the covariance under the physical pairing is $\mathcal C$,
whereas its coordinate covariance is
$\operatorname{Cov}_{\rm count}(Z)=\eta^{-4}\mathcal C$.
The dimensionless derivative vector therefore has exactly
\begin{equation}\label{jt-derivative-covariance}
 \operatorname{Cov}_{\rm count}(\eta^2\nabla Z)
 =\nabla\mathcal C\nabla^T\le K_D I.
\end{equation}
The inequality follows from the proved real order
$\mathcal C\le\Gamma\le G_0$ and the uniform first-derivative
sandwich of $G_0$. Input and output fine spaces have the same
physical scalar measure, so the operator norm of this coordinate
matrix is unchanged on passing to counting Hilbert norms.
The color basis is real trace-orthonormal.

Fixed microscopic blocks contain a fixed number of tests, so
Theorem~\ref{jt-real-theorem} and
\eqref{jt-exponential-square}--\eqref{jt-connected-sum} apply
uniformly over the full homogeneous refinement. Rescaled curls
are fixed local combinations of this derivative vector and
have the same conclusion with fixed norm and incidence factors.
The sampled microscopic H\"older term satisfies
$\eta^{2+\alpha}[\nabla Z]_\alpha
\le C\norm{\eta^2\nabla Z}_\infty$ and can be included at
fixed cost. A unit physical cell contains $M^4$ fine sites,
so treating it as one block here would make $m$ grow with $M$.

For the full macroscopic Gaussian the potential vector instead
has covariance
\begin{equation}\label{jt-potential-covariance}
 \operatorname{Cov}_{\rm count}(\eta Z)=\eta^{-2}\mathcal C.
\end{equation}
Its counting-operator norm is not uniformly bounded over
refinement. For $M\ge3$, at zero coarse momentum a transverse fine alias
$r=e_1$, polarized in direction $2$, lies in the constraints
and has Maxwell eigenvalue $4M^2\sin^2(\pi/M)\le4\pi^2$.
Thus \eqref{jt-potential-covariance} has norm at least
$M^2/(4\pi^2)$. Pointwise potential variance bounds do not
remove this operator obstruction.

For the first bare homogeneous step of a fixed ratio $L$, the
exact physical scaling established in the preceding section
instead gives the joint value/derivative covariance
\begin{equation}\label{jt-fixed-step-covariance}
 \Sigma_{\rm step}=S_L\mathcal C_LS_L^T\le K_LI,
 \qquad S_L=\begin{pmatrix}LI\\\nabla_{1/L}\end{pmatrix}.
\end{equation}
Here $L$ is fixed and the number of retained cells is arbitrary.
The same theorem therefore applies to both components together,
including their cross covariances. The real reduced-Hessian
comparison of that bare step enlarges $K_L$ by at most
$(1-\epsilon)^{-1}$ for chart-coordinate tests. At real
backgrounds it also gives a uniform bound for the separately
identified linearized-potential tests, so the real theorem
applies to those as well.

These applications do not replace the precision in a later
effective step by the bare fixed-$L$ Maxwell form. A later
coarse precision and its separate current terms require their
own covariance identification and uniform bound.

\subsection{A normalized complex quadratic extension}
\label{jt-complex-section}

Whiten the real baseline Gaussian on its nondegenerate slice:
let $\zeta$ be standard real Gaussian in $\mathbb R^N$ and
let $Y=A\zeta$, with $AA^T\le KI$. The observable covariance
may be singular. Suppose $E$ is complex symmetric and
\begin{equation}\label{jt-relative-perturbation}
 E^T=E,\qquad \norm E_{2\to2}\le\epsilon\le\frac14.
\end{equation}
Define the normalized complex Gaussian functional
\begin{equation}\label{jt-complex-measure}
 d\mu_E(\zeta)=\det(I+E)^{1/2}
      e^{-\zeta^TE\zeta/2}\,d\gamma_N(\zeta),
\end{equation}
with the analytic determinant branch equal to one at $E=0$.
Its total mass is one. For a physical quadratic perturbation
$V_X$ of the positive baseline precision $K_X$, condition
\eqref{jt-relative-perturbation} means exactly
$\norm{K_X^{-1/2}V_XK_X^{-1/2}}\le\epsilon$.
An absolute smallness assertion about matrix entries is not a
substitute for this relative bound. At complex $E$ the functional
is not a positive probability law.

\begin{theorem}[Complex Gaussian local-event bound]
\label{jt-complex-theorem}
Under \eqref{jt-incidence} and
\eqref{jt-relative-perturbation}, every set $S$ of $k$ blocks
satisfies
\begin{equation}\label{jt-complex-bound}
 \boxed{\quad
 \left|\mu_E\!\left(\bigcap_{b\in S}\mathsf B_b(R)\right)\right|
 \le\left[2^m\exp\!\left(-\frac{R^2}{6\nu K}\right)\right]^k.
 \quad}
\end{equation}
There is no ambient-$N$ factor. The events use fixed real linear
tests $Y=A\zeta$ and fixed real thresholds.
\end{theorem}

\begin{proof}
All tested coordinates in $S$ span a real row space of dimension
$r\le mk$. Make a real orthogonal change of the whitened
coordinates so that these tests depend only on $x\in\mathbb R^r$;
denote the untested orthogonal coordinates by $z$. Integrating
$z$ exactly gives a normalized complex Gaussian marginal on
$x$ with precision
\begin{equation}\label{jt-marginal-precision}
 B=I+E_{11}-E_{12}(I+E_{22})^{-1}E_{21}.
\end{equation}
Therefore
\begin{equation}\label{jt-marginal-control}
 \norm{B-I}\le\epsilon+\frac{\epsilon^2}{1-\epsilon}
       =\frac{\epsilon}{1-\epsilon}=:s\le\frac13,
 \qquad \operatorname{Re}B\ge(1-s)I=:\beta I,
 \qquad \beta\ge\frac23.
\end{equation}
The real-part assertion is an inequality for the real quadratic
form, equivalently the Hermitian numerical range of the complex
symmetric matrix. Since
$|\det B|\le\norm B^r\le(1+s)^r$, the modulus of the
normalized marginal density is bounded by
\begin{equation}\label{jt-marginal-domination}
 \left(\frac{1+s}{\beta}\right)^{r/2}
                d\gamma_{\beta^{-1}I_r}(x)
 \le 2^{r/2}d\gamma_{\beta^{-1}I_r}(x).
\end{equation}

Under this dominating real Gaussian the tested covariance norm
is at most $K/\beta$. Collect all tests in the selected blocks,
including repetitions. Its dimension is at most $mk$, overlap
at most $\nu$, and covariance norm at most $\nu K/\beta$.
On the all-bad event its squared norm exceeds $kR^2$.
Equation~\eqref{jt-square-mgf}, with
$t=\beta/(4\nu K)$, bounds its probability by
$2^{mk/2}\exp[-\beta kR^2/(4\nu K)]$.
Multiply by $2^{r/2}\le2^{mk/2}$ and use $\beta\ge2/3$.
This proves \eqref{jt-complex-bound}, including degenerate
observable covariance and the case $r=0$.
\end{proof}

The proof also bounds a bounded measurable factor depending
only on the same tested marginal, multiplied by the all-bad
indicator, by its supremum times the right side of
\eqref{jt-complex-bound}. It does not bound the integral of
the modulus of the full $N$-dimensional density over that
event without an $N$-dependent factor. Untested variables must
be integrated out before taking the modulus. An observable
using further coordinates requires enlarging the marginal
and accounting for its dimension.

For fixed real-coordinate events, the normalized integral is
holomorphic in $E$ on the complex operator ball. At each finite
$N$, the real-part lower bound gives Gaussian domination on
compact subballs and the determinant branch is analytic. If
$E$ depends holomorphically on sources in a known neighborhood
remaining inside \eqref{jt-relative-perturbation}, Cauchy on
that neighborhood gives factorial source derivatives of these
local event integrals, retaining the displayed bad-block decay
on a slightly larger parameter domain. This requires a proved
relative perturbation domain; it does not produce one for an
unmatched source Hessian.

In particular an indicator around tests that themselves move
with a complex background is not covered by this holomorphy
argument. For the reduced first-step Hessian, the fixed real
chart-coordinate tests are admissible. Background-dependent
linearized-potential tests are covered by the real covariance
theorem at real backgrounds, but require a separate analytic
cutoff construction at complex backgrounds.

Expansion of a finite product also gives
\begin{equation}\label{jt-complex-product}
 \left|\mu_E\prod_b(1+z_b\mathbf1_{\mathsf B_b(R)})\right|
 \le\prod_b(1+p_{\rm c}(R)z_b),\qquad
 p_{\rm c}(R)=2^m e^{-R^2/(6\nu K)},\qquad z_b\ge0.
\end{equation}
One must not cap $p_{\rm c}$ at one merely because the total
mass is normalized: complex expectations of events are not
probabilities. The same connected-set counting as
\eqref{jt-connected-sum} bounds the sum of absolute local
event expectations whenever its entropy inequality holds.

\subsection{Source match and application boundary}
\label{jt-scope}

For the homogeneous flat derivative application, the action,
coupling and physical pairing are the deposited conventions at
Raw I:464--479, I:12756--12760 and I:15643--15648. The flat
corrected average and its homogeneous composition are matched
at I:24349--24373 and I:47046--47048; the scalar projection is
matched at I:23937--23944. The analytic input actually used is
the established variational $\mathcal C$ and full-$G_0$
first-derivative sandwich. These are checks against the pinned
deposited transcription, not a new original-paper verification.

The real theorem therefore applies to that explicit homogeneous
flat constrained Gaussian, and to the fixed-ratio first bare
step with the covariance identification already proved. The
complex theorem applies under its explicit relative quadratic
hypothesis; for the reduced first bare step this hypothesis is
supplied by the preceding Hessian calculation after shrinking
the background ball. Later effective current-dependent steps
require their own identification.

These joint bounds are stronger than marginal variance or a
single union bound. They still do not integrate nonlinear
interactions, Haar or chart Jacobians, or regions outside the
chart. They do not supply a complete retained-history covariance,
an interacting polymer expansion, or a quantum continuum limit.

% Notebook 17 module. No preamble. Unique label prefix cs-.
% Requires the usual norm macro and jt-real-theorem from the joint-tail module.
\section{The Gaussian at an arbitrary current renormalization scale}
\label{cs-section}

The fixed-ratio bare Gaussian and the Gaussian used after earlier
classical minimizations are different objects. We now extract the
uniform bounds for the latter. They are already contained in
Notebooks 05--07; this section supplies the current-scale probability
normalization and records how Notebook 15 removes the earlier
homogeneous smooth-minimizer hypothesis.

\subsection{The retained coordinates and their exact precision}

Set the current lattice spacing to one. The original spacing is
$M^{-1}$, with $M=L^k$, while the next blocking ratio $L\ge3$ is fixed
and odd. All regulators are compatible periodic tori. Let
\[
 q_M(B)=\inf_{Q_M^{\rm st}A=B}\norm{d_{1/M}A}_{1/M}^2
\]
be the classical flat effective quadratic form. The fine norm has
volume factor $M^{-d}$ and the current norm is counting measure.
Next-step maps are denoted by $Q_L^\flat,T_L,K_L$: respectively the
rooted constraint, averaged root-to-site contour map, and pivot
insertion. These are the fixed-$L$ maps of Notebook 05. The map
$K_L$ inserts a field in $\ker Q_L^\flat$ and is the identity on
the retained nonpivot coordinates. It is not a gauge-selected
fine minimizing extension.

The actual flat exponential-gauge quadratic model on the retained
variables $z$ is
\begin{equation}\label{cs-precision}
 A_{M,0}=K_L^T(q_M+T_L^TT_L)K_L,\qquad
 d\mu_{M,0}(z)=
 (2\pi)^{-n/2}\det(A_{M,0})^{1/2}
 e^{-z^TA_{M,0}z/2}\,dz .
\end{equation}
Here $n$ counts independent real bond/color coordinates once, using
a real color basis orthonormal for the normalized trace. Thus
\begin{equation}\label{cs-counting}
 \operatorname{Cov}_{\rm count}(z)=A_{M,0}^{-1}.
\end{equation}
There is no factor $M^d$ in \eqref{cs-counting}: these are
\emph{current-unit} coordinates, not original fine potential values.
The coupling has already been removed from the Gaussian by the
change of field $B=g_kK_Lz$ at linear order. The nonlinear constraint
correction, Haar density and coordinate Jacobian are kept outside
the chosen Gaussian in \eqref{cs-precision}. In particular, a
coupling-suppressed quadratic Haar term may be part of the
interaction; no vanishing of that term is asserted.

\subsection{Uniform spectral bounds, including all previous depths}

The normalized cochain contraction in Notebook 06 gives
\begin{equation}\label{cs-energy-sandwich}
 \norm{d_1B}_1^2\le q_M(B)\le C_d\norm B_1^2,
 \qquad C_d=2^{d+6}(d-1).
\end{equation}
The compact lift proves the upper bound for odd $M\ge3$; $M=1$
uses the identity lift and satisfies the same displayed bound.
These estimates hold for all current fields, not just a transverse
Fourier sector.

For clarity we record the fixed geometric constants of Notebook 06.
On the cube of $3^d$ neighboring $L$-blocks form the finite matrix
\[
 \mathcal M_L=d_\Omega^Td_\Omega+
 T_{0,\Omega}^TT_{0,\Omega}+
 Q_{{\rm st},\Omega}^TQ_{{\rm st},\Omega},
 \quad \sigma(d,L)=\lambda_{\min}(\mathcal M_L)>0.
\]
Its kernel is trivial: zero curl gives a gradient, the tree rows
make its scalar potential constant on each block, and the averaging
rows equate neighboring block constants. Summing the local
inequality gives the axial gap
$\gamma_{\rm loc}=\sigma(d,L)/3^d$ on
$\ker Q_L^{\rm st}\cap\ker T_0$. Periodic repetitions are counted
with multiplicity, so the same constant applies on short compatible
tori.

Put
\begin{equation}\label{cs-constants}
 \begin{split}
 d_0&=2\sqrt d,\qquad
 t_0=\sqrt{d(L-1)(L^d-1)},\qquad
 r_L=1+2d(L-1)/L,\\
 K_*^2&=1+L^{2d}(2d+2)r_L^2,\\
 \gamma_*&=
 \left((1+d_0t_0)^2/\gamma_{\rm loc}+d_0^2\right)^{-1},
 \qquad B_*=(C_d+t_0^2)K_*^2.
 \end{split}
\end{equation}
The flat gauge bridge of Notebook 05 gives
\[
 q_M(B)+\norm{T_LB}^2\ge\gamma_*\norm B^2
 \quad (Q_L^\flat B=0).
\]
One can see the constraint compatibility directly. First subtract
the rooted gradient $D(T_LB)$ to obtain $B_T$ with
$T_LB_T=Q_L^{\rm st}B_T=0$. Convert it by a block-mean-zero
gradient to the tree representative $B_0$. Their curls agree and
$\norm{B_T}\le(1+d_0t_0)\norm{B_0}$. Apply the axial gap and
weighted Cauchy--Schwarz to
$B=B_T+D(T_LB)$. This argument uses the lower inequality in
\eqref{cs-energy-sandwich}; it does not confuse the rooted and
straight constraints.

Since $z$ is a coordinate restriction of $K_Lz$ and
$\norm{K_L}\le K_*$, the full precision sandwich is
\begin{equation}\label{cs-spectral}
 \gamma_*I\le A_{M,0}\le B_*I,
 \qquad
 B_*^{-1}I\le A_{M,0}^{-1}\le\gamma_*^{-1}I.
\end{equation}
All constants are independent of $k$, $M$ and the torus volume.
There is no need to replace this proved bound by an assumed
transverse ellipticity statement.

\subsection{Weighted locality and its scope}

Group the retained coordinates by their fixed next-step $L$-cells.
For a block matrix $F$, use the weighted Schur norm
\[
 \norm F_\kappa=
 \max\left\{
 \sup_i\sum_j e^{\kappa|i-j|_1}\norm{F_{ij}},
 \sup_j\sum_i e^{\kappa|i-j|_1}\norm{F_{ij}}
 \right\},
\]
with periodic distance on finite regulators. Notebook 07,
its flat-decay theorem and section on periodic gluing and the reduced inverse,
prove constants $\kappa_*>0,C_*<\infty$, depending only on fixed
geometry, such that
\begin{equation}\label{cs-weighted}
 \norm{A_{M,0}}_{\kappa_*}
 +\norm{A_{M,0}^{-1}}_{\kappa_*}\le C_*.
\end{equation}
The proof there first establishes a uniform holomorphic Fourier
strip for $q_M$, then uses fixed-range $K_L,T_L$, the real gap
\eqref{cs-spectral}, and a smaller strip for the inverse.
Exact periodization gives the same conclusion on compatible tori.
This additional locality does not follow from a spectral sandwich
alone. It concerns the flat retained-coordinate matrix; it is not
strict locality in a varying background.

\subsection{Joint value-and-difference tails on the current lattice}

Let $\nabla_1$ contain all forward unit differences of every bond
component, and define
\begin{equation}\label{cs-tested-field}
 Y=\begin{pmatrix}K_Lz\\\nabla_1K_Lz\end{pmatrix}.
\end{equation}
The elementary bound
$\norm{\nabla_1}\le2\sqrt d$ and
\eqref{cs-spectral} imply
\begin{equation}\label{cs-tested-covariance}
 0\le\operatorname{Cov}_{\rm count}(Y)\le
 K_{\rm cur}I,
 \qquad K_{\rm cur}=(1+4d)K_*^2/\gamma_*.
\end{equation}
Because $K_L$ and the difference stencil have fixed range,
\eqref{cs-weighted} also supplies exponential block decay for this
tested covariance, with a fixed loss in its constant.

If each tested cell uses at most $m$ scalar coordinates and every
coordinate is tested in at most $\nu$ cells, the real joint-tail
theorem~\ref{jt-real-theorem} gives
\begin{equation}\label{cs-joint-tail}
 \mathbb P\left(
 \max_{i\in I_b}|Y_i|>R\text{ for every }b\in S\right)
 \le
 \left[m\sqrt2\exp\left\{-\frac{R^2}{4\nu K_{\rm cur}}\right\}
 \right]^{|S|}.
\end{equation}
Fixed current-scale cells, including fixed collars for derivatives,
have $m,\nu$ independent of $M$. The Gaussian is correlated;
independence is neither asserted nor needed. For $g_kY$, replace
$R$ by $R/g_k$. If $R(g)=c(\log g^{-2})^p$ with $p>1/2$, the
bracket at radius $R(g)$ is smaller than any fixed power of $g$
below a sufficiently small ceiling.

To restore the physical current spacing $a$, write
$Z_{\rm phys}(x)=a^{-1}(K_Lz)(x/a)$ in four dimensions. Then
\[
 (aZ_{\rm phys},a^2\nabla_aZ_{\rm phys})=Y.
\]
This is the current-step normalization. It differs from treating
all original fine variables between $M^{-1}$ and one as a single
fluctuation. It also differs from the bare fixed-ratio formula
$(LI,\nabla_{1/L})C_L(LI,\nabla_{1/L})^T$, whose coordinates use
next-coarse spacing one. Both statements are consistent after their
scales and Gaussian models are specified.

\subsection{What is now identified with the deposited source}

Notebook 07's source section records the chain from the deposited
small-field step to \eqref{cs-precision}: Raw II:73865--73950
straightens the constraint, rescales by $g_k$, and eliminates the
pivots; Raw I:46960--47000 fixes the pivot coefficient convention.
At the trivial field the action Hessian is $q_M$, the exponential
gauge Hessian is $T_L^TT_L$, and the action gradient vanishes.
Consequently the quadratic start of the constraint-straightening
correction does not change this Hessian. Haar and Jacobian factors
remain in the interaction.

The smooth-minimizer premise in that earlier identification is now
supplied, in the homogeneous setting, by Notebook 15's theorem on
transfer to the ordinary homogeneous averaging constraint. It
constructs an analytic local branch for the ordinary iterated group
average, with real action minimality in a fixed small global
potential class. At each finite regulator the ordinary constraint
is an analytic submersion at the identity. Its derivative is onto,
as follows from the fixed-step pivot right inverses and their
composition. Thus every tangent-feasible field is the tangent to
a feasible curve. Comparing the local minimum with these curves
shows that its second derivative is the quadratic variational
infimum for the rooted derivative. Notebook 06's explicit
gradient adjustment identifies that infimum with the displayed
$Q_M^{\rm st}$ infimum. The gradient adjustment preserves energy;
it does not identify the two averaging maps as maps on potentials.

Accordingly \eqref{cs-precision} is the homogeneous flat source
Gaussian under the deposited ordinary-average and normalization
conventions already matched in those notebooks. The nonlinear
branch need not equal an independently chosen axial representative:
the locally minimizing action and its flat Hessian agree modulo
gauge. This does not re-certify every original journal convention,
extend the result to mixed histories or holes, or identify a
nonflat independent-current analytic family.

\subsection{A distinct fixed-slice perturbation consequence}

There is also a useful quadratic model without pivot coordinates.
Let $X_L=\ker Q_L^{\rm st}\cap\ker(R_LD_1^T)$ be the fixed flat
source-gauge slice for the unit-to-$L$ step, and let $J_L$ be its
isometric inclusion. The fixed-$L$ Maxwell covariance proved in
Notebooks 09 and 12 implies
\[
 J_L^Td_1^Td_1J_L\ge c_LI
\]
for some $c_L>0$, uniformly in compatible volume. For example one
may take $c_L=(L^2\norm{C_L})^{-1}$, where $C_L$ uses next-coarse
spacing one. Hence
\begin{equation}\label{cs-slice-gap}
 c_LI\le J_L^Tq_MJ_L\le C_dI.
\end{equation}
This gives the same kind of current-scale value-and-difference
tail bound on $X_L$, uniformly in $M$.

Notebook 16 proves that the ordinary coarse classical Hessian
satisfies $D^2V_M(B)=q_M+O(\norm B_\infty)$ in weighted row and
column norms. Its compression to this fixed slice therefore has a
relative perturbation bounded by $C\norm B_\infty/c_L$. This is
a valid input to the abstract complex relative-quadratic tail
theorem on a sufficiently small common background ball. It is an
effective classical Gaussian family in a fixed linear slice.
It is not automatically the source's nonlinear fiber Gaussian:
the latter also contains chart-curvature/current terms, exponential
gauge fixing, and an exact cutoff-coordinate prescription.
Notebook 19 now performs that additional physical-branch
identification, keeping both chart-curvature terms and the
actual absolute gauge function.
Nor is weighted locality of this compressed inverse inferred by
conjugating with weights which do not preserve $X_L$; the weighted
inverse assertion \eqref{cs-weighted} concerns the separately
proved pivot-coordinate family.

\section{The next integration and iteration}
\label{gs17-next}

For a logarithmic cutoff $R(g)=(\log(1/g))^p$ with $p>1/2$,
the Gaussian bad-block factor is smaller than any fixed power of
$g$, while $gR(g)^3$ tends to zero. These two estimates are
compatible: the same cutoff can suppress bad blocks and make
analytic cubic action remainders small. Jacobian terms that begin
linearly require their own local bounds; they are not cubic merely
because the action remainder is cubic.

The quantum step requires independent fixed-dimensional integration
coordinates, their normalized covariance, and localized estimates
for the complete interaction. The flat pivot covariance above
supplies such coordinates on the homogeneous regulator. The
singular value/gradient vector is useful for tests, but is not an
independent cell coordinate system for a Gaussian density expansion.
For complex continuation the tests used here remain fixed on the
real integration variables.

Notebook 18 gives a convergent cutoff expansion with overlapping
interaction supports under an explicit weighted coefficient bound.
Notebook 19 identifies the nonflat physical Gaussian and bounds
the local straightening, Haar and gauge factors. Notebook 20
verifies the complete coefficient and precision hypotheses for
the bare Wilson small-field step, with all normalization terms.
Notebook 21 supplies strict localization of the minimized
classical action uniformly in earlier depth and closes one
integration of that reference action. Later quantum steps
require transport to the next coarse background and control
of the accumulated interactions. Bad regions require
regional constructions and their normalization estimates. Iterating
those bounds must also control the running couplings and preserve
the observable hierarchy needed by the continuum construction.

The output sought from the iteration remains a common nontrivial
gauge-invariant continuum hierarchy with the Euclidean covariance,
positivity, regularity and Yang--Mills identification specified in
Notebook 01. The gap argument is a further target after that
existence construction is closed.

\begin{thebibliography}{9}
\bibitem{prior17}
\emph{Toward a four-dimensional Yang--Mills construction},
research Notebooks 5--16 in this workspace, 16 October 2026.
The flat reduced covariance, constrained Maxwell inverse,
nonlinear minimum and classical response theorems are used here.
\bibitem{source17}
E.~Sussman and collaborators,
\href{https://github.com/ethanwsussman/Yang--Mills-LLM-experiment}{Yang--Mills-LLM-experiment},
deposited raw accounts at commit
\texttt{e392ef45013a581664bab76b90867170b30fe10e}.
The Yang--Mills files are unchanged at checked commit
\texttt{785a68fdeef78a1da854c4807fe1f17b597d05e7}.
\bibitem{balaban17}
T.~Ba\l aban, \emph{Renormalization group approach to lattice gauge
field theories I}, Commun. Math. Phys. \textbf{109} (1987), 249--301,
\href{https://scispace.com/pdf/renormalization-group-approach-to-lattice-gauge-field-3ygai0vxt1.pdf}{original article copy}.
The source-coordinate match is discussed explicitly above.
\end{thebibliography}
\end{document}
